If the function $U(x)$ has a singularity at some real $x$, the reflection
coefficient is governed mainly by the field near that point. To calculate it,
perturbation theory may be applied formally without requiring its condition of
applicability to hold for every $x$; it is enough that the quasiclassical
condition be satisfied. We then arrive at formula (1) of Problem~2, with the
sole difference that the value of the function $p(x)$ at $x=x_0$ must appear
there in place of the incident particle's momentum.
